\section{对比学习：更通用的表示目标}

\subsection{从“同一对象的不同视图”出发}

如果说图像变换类预文本任务是在问“这张图被怎么变了”，那么对比学习问的是“哪些样本应该被拉近，哪些样本应该被推远”。

\begin{figure}[H]
\centering
\includegraphics[width=0.95\textwidth]{slides-images/slide-066.jpg}
\caption{对比学习的直觉：同一对象的不同视图相互吸引，不同对象彼此排斥\protect\footnotemark}
\end{figure}
\footnotetext{Slides 第66页。}

\begin{knowledgebox}{对比学习的目标}
\begin{itemize}
    \item 正样本：同一对象的不同增强视图；
    \item 负样本：其他对象；
    \item 目标：正样本在表示空间里更近，负样本更远。
\end{itemize}
\end{knowledgebox}

\subsection{InfoNCE 形式化}

给定参考样本 $x$、正样本 $x^+$、负样本 $x^-$，令编码器为 $f(\cdot)$，打分函数为 $s(\cdot,\cdot)$。希望
\[
s(f(x), f(x^+)) \gg s(f(x), f(x^-)).
\]
对应的常见损失就是 \textbf{InfoNCE}：
\[
\mathcal{L}
= - \log \frac{\exp(s(f(x), f(x^+))/\tau)}
{\exp(s(f(x), f(x^+))/\tau) + \sum_{i=1}^{N-1}\exp(s(f(x), f(x_i^-))/\tau)}.
\]

\begin{figure}[H]
\centering
\includegraphics[width=0.95\textwidth]{slides-images/slide-074.jpg}
\caption{InfoNCE 可以看作一个 N-way softmax 分类问题\protect\footnotemark}
\end{figure}
\footnotetext{Slides 第74页。}

从形式上看，它和多分类交叉熵几乎一样：模型要在一堆候选里把“正确的那个正样本”挑出来。

\begin{importantbox}{为什么 InfoNCE 有效}
它等价于“把正样本当作正确类别、把所有负样本当作其他类别”的分类问题，因此特别适合端到端优化。
\end{importantbox}

\subsection{为什么需要大 batch}

InfoNCE 的一个现实问题是：负样本越多，通常效果越好。因此很多方法都依赖大 batch 或者大字典/队列来提供更多负样本。

\begin{knowledgebox}{负样本越多通常越好}
更多负样本通常会带来更紧的互信息下界，也让模型更难走捷径，因此学习到的表示往往更强。但代价是显存和训练复杂度都会上升。
\end{knowledgebox}

\subsection{SimCLR}

SimCLR 是一个非常简洁但很强的实例级对比学习框架。

\begin{figure}[H]
\centering
\includegraphics[width=0.95\textwidth]{slides-images/slide-076.jpg}
\caption{SimCLR：数据增强、projection head 和 cosine similarity\protect\footnotemark}
\end{figure}
\footnotetext{Slides 第76页。}

它的关键点有三个：
\begin{itemize}
    \item 用强数据增强生成正样本对；
    \item 在 encoder 后加一个 projection head $g(\cdot)$，只在投影空间做对比学习；
    \item 用一个大 batch，把 batch 里其他样本都当作负样本。
\end{itemize}

\begin{importantbox}{为什么要 projection head}
对比损失有时会丢掉一些对下游任务有用的信息。projection head 允许模型把“对比学习所需的表示”和“下游任务最有用的表示”适当分开。
\end{importantbox}

\subsection{SimCLR 的训练细节}

SimCLR 的训练流程可以理解为：
\begin{enumerate}
    \item 对每张图做两次随机增强；
    \item 得到 2N 个视图；
    \item 每个视图把另一个增强视图当正样本，其余 2N-2 个当负样本；
    \item 用 InfoNCE 逐个计算损失并平均。
\end{enumerate}

\begin{figure}[H]
\centering
\includegraphics[width=0.95\textwidth]{slides-images/slide-081.jpg}
\caption{SimCLR 的 mini-batch 结构：每对增强视图构成一个正样本对\protect\footnotemark}
\end{figure}
\footnotetext{Slides 第81页。}

SimCLR 的另一个重要结论是：\textbf{强增强和大 batch 都非常重要}。没有足够多负样本，效果会明显下降。

\subsection{MoCo 与 MoCo v2}

SimCLR 很强，但对 batch size 和显存要求很高。MoCo 的思路是把负样本从当前 batch 里“解耦”出来，使用一个队列来持续维护大量 key。

\begin{figure}[H]
\centering
\includegraphics[width=0.95\textwidth]{slides-images/slide-088.jpg}
\caption{MoCo：用 queue 维护大量负样本，并通过 momentum encoder 更新 key 分支\protect\footnotemark}
\end{figure}
\footnotetext{Slides 第88页。}

MoCo 的要点是：
\begin{itemize}
    \item query 分支负责反向传播；
    \item key 分支不直接回传梯度，而是用 momentum update 慢慢更新；
    \item 队列里的旧 key 作为负样本，解决了 batch 受限的问题。
\end{itemize}

\begin{importantbox}{MoCo 的工程价值}
MoCo 把“负样本规模”和“batch size”分开了，因此比 SimCLR 更省显存，也更适合普通硬件。
\end{importantbox}

\begin{figure}[H]
\centering
\includegraphics[width=0.95\textwidth]{slides-images/slide-090.jpg}
\caption{MoCo v2：把 SimCLR 的强增强和 projection head 引入 MoCo\protect\footnotemark}
\end{figure}
\footnotetext{Slides 第90页。}

MoCo v2 本质上是一个混合体：既保留 MoCo 的队列和 momentum encoder，又吸收了 SimCLR 的非线性投影头和更强的数据增强。

\subsection{序列级对比学习：CPC}

前面的 SimCLR 和 MoCo 都是 \textbf{instance-level} 对比学习。CPC 则把对比目标放在时序上。

\begin{figure}[H]
\centering
\includegraphics[width=0.95\textwidth]{slides-images/slide-094.jpg}
\caption{Instance-level 与 sequence-level 对比学习的区别\protect\footnotemark}
\end{figure}
\footnotetext{Slides 第94页。}

CPC 的思路是：先把一个序列编码成局部向量 $z_t$，再把前半段上下文汇总成 context code $c_t$，最后用 InfoNCE 让 $c_t$ 更容易预测未来的 $z_{t+k}$，而不是错误的 future code。

\begin{figure}[H]
\centering
\includegraphics[width=0.95\textwidth]{slides-images/slide-098.jpg}
\caption{CPC 的基本流程：编码局部表示，汇总上下文，再预测未来\protect\footnotemark}
\end{figure}
\footnotetext{Slides 第98页。}

\begin{knowledgebox}{CPC 的适用场景}
\begin{itemize}
    \item 音频序列；
    \item 视频帧序列；
    \item 一些可以按顺序组织的视觉任务。
\end{itemize}
\end{knowledgebox}

不过在图像表示学习上，CPC 往往不如更直接的实例级方法强。这也是为什么后来的主流更多转向 SimCLR、MoCo、MAE 和 DINO 这类方法。

\subsection{DINO}

DINO 不是严格意义上的 contrastive learning，但它借用了“student/teacher、视图一致性”的思想。

\begin{figure}[H]
\centering
\includegraphics[width=0.95\textwidth]{slides-images/slide-103.jpg}
\caption{DINO：Self-Distillation with No Labels\protect\footnotemark}
\end{figure}
\footnotetext{Slides 第103页。}

它的主要特点是：
\begin{itemize}
    \item 不显式构造传统的正负样本对；
    \item 通过 teacher-student 机制让不同视图输出一致；
    \item 常与 ViT 结合，学习到的特征往往非常适合可视化和下游迁移。
\end{itemize}

\begin{importantbox}{DINO 的位置}
DINO 继承了对比学习的“让不同视图保持一致”的思想，但它更接近自蒸馏，而不是标准的负样本对比框架。
\end{importantbox}

%% ========== 总结 ==========

